% =============================================================================
% 05_data_and_metadata.tex -- Data and Metadata (~4-5 frames)
% =============================================================================

\section{Data and Metadata}

% ---- Frame 1 ----
\begin{frame}{The BTC Dataset (OpenNeuro ds002080)}

  \textbf{Brain Tumor Connectomics} -- Aerts et al.\
  \cite{Aerts2018,Aerts2022}

  \bigskip

  \begin{itemize}
    \item 25 glioma patients, pre- and post-operative mpMRI
    \item T1w, multi-shell DWI, resting-state BOLD per session
    \item We use \textbf{1 subject} $\times$ \textbf{2 timepoints}:
  \end{itemize}

  \bigskip

  \begin{tabular}{llll}
    \toprule
    \textbf{Subject} & \textbf{Session} & \textbf{Tumor} & \textbf{Modalities} \\
    \midrule
    PAT16 & pre-op  & LGG, left hemisphere & T1w, DWI (AP+PA), BOLD \\
    PAT16 & post-op &  (resected)          & T1w, DWI (AP+PA), BOLD \\
    \bottomrule
  \end{tabular}

  \medskip

  \textbf{Phenotypic data} (per session):
  psychosocial questionnaires (STAI, PSWQ, ERQ),
  CANTAB cognitive battery (MOT, RTI, RVP, SOC, SSP)
  $\to$ pre-/post-op \alert{cognitive outcome} tracking.

\end{frame}

% ---- Frame 2 ----
\begin{frame}[fragile]{BIDS Organization}

  Data organized according to the \textbf{Brain Imaging Data Structure}:

  \bigskip

  \begin{columns}[T]
    \column{0.5\textwidth}
    \begin{lstlisting}
data/BIDS_preop_postop/
  participants.tsv
  sub-PAT16/
    sub-PAT16_sessions.tsv
    ses-preop/
      anat/  *_T1w.nii.gz
      dwi/   *_acq-AP_dwi.nii.gz
      func/  *_task-rest_bold...
    ses-postop/
      anat/  ...
      dwi/   ...
      func/  ...
    \end{lstlisting}

    \column{0.45\textwidth}
    \begin{itemize}
      \item Subject $\times$ session hierarchy
      \item AP + PA DWI pairs $\to$ robust distortion correction
      \item JSON sidecars with acquisition parameters
      \item \texttt{*\_sessions.tsv}: 48 cols of clinical \& cognitive data
      \item \texttt{download\_data.py} retrieves from OpenNeuro
    \end{itemize}
  \end{columns}

\end{frame}

% ---- Frame 3 ----
\begin{frame}{Acquisition Parameters}

  \begin{tabular}{llll}
    \toprule
    \textbf{Modality} & \textbf{Sequence} & \textbf{Resolution} & \textbf{Key parameters} \\
    \midrule
    T1w  & MPRAGE & 1\,mm iso & $176 \times 256 \times 256$ \\
    DWI  & EPI    & 2.5\,mm iso & $b = 700,\, 1000,\, 2000$\,s/mm$^2$ \\
    BOLD & EPI    & 3\,mm iso & TR $= 1.7$\,s, resting state \\
    \bottomrule
  \end{tabular}

  \bigskip

  \textbf{Key differences from the healthy-volunteer project:}
  \begin{itemize}
    \item Multi-shell DWI (3 shells vs.\ 4 in healthy)
    \item AP + PA DWI pairs available $\to$ topup-based SDC
    \item Pathological anatomy in all sessions
    \item Two timepoints per subject $\to$ longitudinal analysis
  \end{itemize}

\end{frame}

% ---- Frame 4 ----
\begin{frame}{What the Data Gives Us for Brain Mechanics}

  \begin{center}
  \begin{tikzpicture}[
    box/.style={draw, rounded corners, minimum width=4cm,
                minimum height=0.8cm, align=center, font=\small},
    arr/.style={-{Stealth[length=3mm]}, thick},
    node distance=1.2cm
  ]
    \node[box] (raw) {Raw mpMRI\\(T1w + DWI + BOLD)};
    \node[box, below=of raw] (deriv) {Co-registered derivatives\\in T1w space};
    \node[box, below left=1.2cm and -0.5cm of deriv]
      (pre) {Pre-op model};
    \node[box, below right=1.2cm and -0.5cm of deriv]
      (post) {Post-op model};
    \node[box, below=1.2cm of deriv] (delta) {\alert{$\Delta$ = empirical deformation}};

    \draw[arr] (raw) -- (deriv);
    \draw[arr] (deriv) -- (pre);
    \draw[arr] (deriv) -- (post);
    \draw[arr] (pre) -- (delta);
    \draw[arr] (post) -- (delta);
  \end{tikzpicture}
  \end{center}

  The longitudinal design provides what most brain mechanics studies
  lack: \alert{observed displacement data} for validation.

\end{frame}
